\section{Limitations and Future Work}

\textbf{Limitations.}~~(1) PCGT is \emph{transductive}: METIS partitions are computed on the full graph (including test nodes) before training. This is consistent with the standard evaluation protocol---all baselines (GCN, GAT, SGFormer) also access full graph structure during message passing---but PCGT cannot directly classify nodes unseen at partitioning time without re-partitioning or a fallback strategy. (2) The current implementation loops over partitions in Python; a padded-tensor masked attention would be faster on GPU (same algorithmic complexity, purely an engineering issue). (3) PCGT runs on ogbn-arxiv (169K nodes) without trouble, but billion-node graphs (ogbn-papers100M) would need partition sampling. (4) The number of partitions $K$ is tuned per dataset; an adaptive rule for setting $K$ would be more practical. (5) On graphs with near-random homophily (Deezer, $h \approx 0.53$), the gain over SGFormer is only +0.1\%---partitioning adds little when the edge labels are close to random.

\textbf{Future work.}~~(1)~Scaling to billion-node graphs (ogbn-papers100m) via partition sampling---process a random subset of partitions per forward pass with EMA-updated representatives. (2)~Batched partition attention using padded tensors and attention masks to eliminate the Python for-loop. (3)~Extension to link prediction and graph classification tasks. (4)~Theoretical analysis connecting partition attention to spectral properties of the graph Laplacian. (5)~Visualization of learned seed representations to interpret what ``aspects'' each seed captures across different partitions.

\textbf{Broader impact.}~~PCGT is a general-purpose node classification method evaluated on standard academic benchmarks.  It does not introduce new data collection or surveillance capabilities.  Like any graph-based model deployed in downstream applications (e.g., social-network recommendation, fraud detection), practitioners should audit predictions for fairness across demographic groups, since biases present in the input graph or features can propagate through learned attention.  The METIS partitioning step is deterministic given the adjacency matrix, so results are fully reproducible.  We are not aware of direct negative societal consequences specific to the partition-conditioned attention mechanism itself.
